\section{Discussion}

MouseBrainBench has been designed to be useful in a strict scientific sense. It
does not only search for positive results, but also identifies the claims that
must be rejected. This behaviour is essential in the current context, where
digital-brain terminology can easily lead to interpretations that are stronger
than the evidence available.

The most relevant positive result of the proposal is local and specific. In
MICRONS CAVE data, synaptically connected directed pairs are more similar in
readout location than distance- and degree-matched non-connected pairs. This
result is replicated in a discovery cohort and two independent hold-out cohorts.
Therefore, it supports a local observational structure-function association in
mouse visual cortex.

However, the result does not support a causal claim. The analysis is based on
observational anatomical and functional data, and there is no perturbation
experiment proving that the detected synaptic connections are necessary or
sufficient for the functional similarity. This point is not a weakness of the
paper if it is stated correctly. On the contrary, it is one of the reasons for
using a claim-audit framework.

The negative and predictive layers also provide relevant information. Allen
Visual Behavior Neuropixels shows that reproducibility alone is insufficient for
mechanistic identifiability. Sensorium static shows that prediction and
reliability can be evaluated in a modern benchmark without converting these
metrics into causal evidence. Dynamic Sensorium shows that temporal predictive
models can improve over simpler baselines while still failing the reliability
and mechanistic gates.

Taken together, these results suggest that partial digital models of the mouse
brain can be scientifically useful if their claims are correctly bounded. The
contribution of MouseBrainBench is precisely this critical boundary: it provides
a reproducible way to decide whether a result is predictive, reproducible,
structurally controlled, locally structure-function positive, or
mechanistically identifiable.
